\documentclass{article}
\usepackage{graphicx} % Required for inserting images
\usepackage[margin=0.5in]{geometry}


\title{CMPT 310 Milestone 1}
\author{\textbf{Group 14:} Finley Thomson, Monem Moheq, Srinivas Suggu, Jarod Hwang}
\date{July 2026}

\begin{document}

\maketitle

\section*{Project Summary}

\quad Our project aims to be able to classify the return on investment of a potential infill townhouse/low-rise real estate development site in Calgary as either high, medium, or low based on neighbourhood, proximity to downtown, proximity to transit stops, and the neighbourhood median income. The labelling for this will be done by taking:

\begin{equation}
    \frac{FinalValue - (InitialValue + Development Costs)}{InitialValue + Development Costs}
\end{equation}
    
We will accomplish this by using an ordinal logistic classifier and a random forest classifier, that will then decide the ROI classification using a soft voting system.

\section*{Accomplishments}

\quad Thus far, we filtered through and gathered data on 225 different real-world examples of infill townhouse/low-rise development projects started and completed between 2017 and 2023. These data points include initial value of the parcel prior to development, development cost, stabilized value following development, longitude, latitude, neighbourhood, proximity to downtown, proximity to the nearest transit stop, and neighbourhood median income. 

In addition to this real world data set, we have formed a larger, artificial dataset formed by using the distribution of values in the categories mentioned above to randomly generate several thousand data points. Although these are not necessarily representative of the real-world (as thousands of extra sites would affect valuation, thus rendering these completely inaccurate) and are not on par with the real world data, we hope to use this supplementary data set in a separate case trial to demonstrate the model's capabilities. 


\section*{Challenges}

\quad The most notable and the major roadblock we faced was synthesizing various municipal data sets into one aggregate set which we can then use to train our model. One more minor roadblock was sorting the building permit data from individual buildings into separate projects, solved by using the DBSCAN algorithm from \texttt{scikit-learn} to classify buildings within 70 meters of the mean as a part of the same projects. A more major roadblock was synthesizing historical assessment data with the building permit data. This was difficult as we had to keep track of the project addresses since, during the development cycle, the parcels are subdivided, meaning that keeping track of which assessment belonged to which address was difficult to match. Furthermore, often multiple of the same addresses were given for any one building permit (one for the parcel of land and one for the building), thus meaning double counting had to be avoided. Ultimately, we managed to sort through the data and assign the correct initial and final assessment values. 


\section*{Project Changes}

\quad Not much has changed about the overall aim of our project, however, the scope has been narrowed somewhat to townhouse/low-rise multifamily buildings in particular, though the idea remains the same. Beyond this, our milestone 1 did have to shift to gathering and processing data, rather than having a rudimentary form of the model running, given the challenges faced when dealing with the municipal datasets.


\end{document}
